% --- Archon physics formalization source begin ---
% archon:physics
% archon:covers PhyXMiniProblems/problem_phyx_mini_0281.lean
% archon:source-report reports/phyx_mini/problem_phyx_mini_0281.source.json

\chapter{Physics problem phyx\_mini\_0281}
\label{ch:PhyXMiniProblems_problem_phyx_mini_0281}

\paragraph{Problem source.}
\#\# Physical scenario

Figure shows that if we hang a block on the end of a spring with spring constant \$k\$, the spring is stretched by distance \$h = 2.0\$ cm. If we pull down on the block a short distance and then release it, it oscillates vertically with a certain frequency.

\#\# Question

What length must a simple pendulum have to swing with that frequency?

\#\# Answer choices

- A: A: 1.60 cm

- B: B: 1.80 cm

- C: C: 2.20 cm

- D: D: 2.00 cm

\#\# Figure caption (auxiliary; use the image as primary evidence)

The image contains three sections:

1. **Left Section:**
   - A spring is hanging from a brick surface.
   - The spring is labeled with "k," possibly indicating its spring constant.
   - Below the spring, there is an arrow pointing downward labeled "h," which might indicate distance or displacement.
   - The spring is freely hanging without any object attached to its end.

2. **Middle Section:**
   - Similar to the left section but with a significant difference.
   - The same spring labeled "k" is shown stretched, with a solid square block attached to its end, indicating it is under tension or bearing weight.

3. **Right Section:**
   - Depicts a pendulum setup.
   - A string is hanging from a brick surface, through a pulley, with a spherical object at the end (likely a pendulum's bob).
   - The string's path suggests the pendulum can swing freely.

The brick surfaces in all sections suggest they serve as the fixed support for the attached components.

\#\# Dataset metadata

Resonance and Harmonics; Multi-Formula Reasoning, Physical Model Grounding Reasoning

\paragraph{Recorded answer/context.}
D: D: 2.00 cm

\paragraph{Figure/image path.}
/root/proposal\_for\_physic/science-mango/phyx\_mini\_run/phyx\_data/test\_image/281.png

\paragraph{Formalization target.}
create a compiling Lean file with sorry bodies at `PhyXMiniProblems/problem\_phyx\_mini\_0281.lean`. The Lean declarations must preserve the physical quantities, dimensions or dimensional roles, figure labels, governing-law hypotheses, and final relation expressed by this problem.
Use Mathlib/Physlib names found through LeanExplore where available. If a physics API is missing, introduce faithful local abstractions rather than scalar placeholder aliases.

\begin{theorem}[Physics formalization target]
\label{thm:physics:phyx_mini_0281:target}
\lean{PhyXMiniProblems.ProblemPhyXMini0281.problem_phyx_mini_0281}
\uses{def:physics:phyx-mini-0281:phyxminiproblems-problemphyxmini0281-lengthincentimeters, def:physics:phyx-mini-0281:phyxminiproblems-problemphyxmini0281-matchesprimaryfigure, def:physics:phyx-mini-0281:phyxminiproblems-problemphyxmini0281-springpendulumsetup, def:physics:phyx-mini-0281:phyxminiproblems-problemphyxmini0281-matchesproblemdescription, def:physics:phyx-mini-0281:phyxminiproblems-problemphyxmini0281-hasphysicalparameters, def:physics:phyx-mini-0281:phyxminiproblems-problemphyxmini0281-satisfiesspringandpendulumlaws, def:physics:phyx-mini-0281:phyxminiproblems-problemphyxmini0281-pendulummatchesspringfrequency, def:physics:phyx-mini-0281:phyxminiproblems-problemphyxmini0281-recordedanswerchoice, def:physics:phyx-mini-0281:phyxminiproblems-problemphyxmini0281-matchesanswerchoice}
The assigned autoformalize agent should translate this physics problem into Lean declarations in the covered file, with theorem and lemma proof bodies written as `by sorry`.
\end{theorem}
\begin{proof}
This is an autoformalization task, not a proof task. Produce statements that can later be proved by the physics prover without weakening the physical model.
\end{proof}

% --- Archon declaration topology begin ---
\section{Declaration topology}
\begin{definition}[Mass Quantity]
  \label{def:physics:phyx-mini-0281:phyxminiproblems-problemphyxmini0281-massquantity}
  \lean{PhyXMiniProblems.ProblemPhyXMini0281.MassQuantity}
  A nonnegative physical mass
\end{definition}
\begin{proof}
  This is the declared model definition of Mass Quantity.
\end{proof}
\begin{definition}[Length Quantity]
  \label{def:physics:phyx-mini-0281:phyxminiproblems-problemphyxmini0281-lengthquantity}
  \lean{PhyXMiniProblems.ProblemPhyXMini0281.LengthQuantity}
  A nonnegative physical length
\end{definition}
\begin{proof}
  This is the declared model definition of Length Quantity.
\end{proof}
\begin{definition}[Velocity Quantity]
  \label{def:physics:phyx-mini-0281:phyxminiproblems-problemphyxmini0281-velocityquantity}
  \lean{PhyXMiniProblems.ProblemPhyXMini0281.VelocityQuantity}
  A nonnegative physical speed or one-dimensional velocity magnitude
\end{definition}
\begin{proof}
  This is the declared model definition of Velocity Quantity.
\end{proof}
\begin{definition}[Acceleration Quantity]
  \label{def:physics:phyx-mini-0281:phyxminiproblems-problemphyxmini0281-accelerationquantity}
  \lean{PhyXMiniProblems.ProblemPhyXMini0281.AccelerationQuantity}
  A nonnegative acceleration magnitude, with dimension L T⁻²
\end{definition}
\begin{proof}
  This is the declared model definition of Acceleration Quantity.
\end{proof}
\begin{definition}[Spring Stiffness Quantity]
  \label{def:physics:phyx-mini-0281:phyxminiproblems-problemphyxmini0281-springstiffnessquantity}
  \lean{PhyXMiniProblems.ProblemPhyXMini0281.SpringStiffnessQuantity}
  Spring stiffness, equivalently force per length, with dimension M T⁻²
\end{definition}
\begin{proof}
  This is the declared model definition of Spring Stiffness Quantity.
\end{proof}
\begin{definition}[Frequency Quantity]
  \label{def:physics:phyx-mini-0281:phyxminiproblems-problemphyxmini0281-frequencyquantity}
  \lean{PhyXMiniProblems.ProblemPhyXMini0281.FrequencyQuantity}
  Cyclic frequency, with dimension inverse time
\end{definition}
\begin{proof}
  This is the declared model definition of Frequency Quantity.
\end{proof}
\begin{definition}[Angular Frequency Quantity]
  \label{def:physics:phyx-mini-0281:phyxminiproblems-problemphyxmini0281-angularfrequencyquantity}
  \lean{PhyXMiniProblems.ProblemPhyXMini0281.AngularFrequencyQuantity}
  Angular frequency; radians are dimensionless
\end{definition}
\begin{proof}
  This is the declared model definition of Angular Frequency Quantity.
\end{proof}
\begin{definition}[length Readout]
  \label{def:physics:phyx-mini-0281:phyxminiproblems-problemphyxmini0281-lengthreadout}
  \lean{PhyXMiniProblems.ProblemPhyXMini0281.lengthReadout}
  \uses{def:physics:phyx-mini-0281:phyxminiproblems-problemphyxmini0281-lengthquantity}
  Read a nonnegative physical length in a selected length unit
\end{definition}
\begin{proof}
  This is the declared model definition of length Readout.
\end{proof}
\begin{definition}[length In Meters]
  \label{def:physics:phyx-mini-0281:phyxminiproblems-problemphyxmini0281-lengthinmeters}
  \lean{PhyXMiniProblems.ProblemPhyXMini0281.lengthInMeters}
  \uses{def:physics:phyx-mini-0281:phyxminiproblems-problemphyxmini0281-lengthquantity, def:physics:phyx-mini-0281:phyxminiproblems-problemphyxmini0281-lengthreadout}
  Metre readout of a physical length
\end{definition}
\begin{proof}
  This is the declared model definition of length In Meters.
\end{proof}
\begin{definition}[length In Centimeters]
  \label{def:physics:phyx-mini-0281:phyxminiproblems-problemphyxmini0281-lengthincentimeters}
  \lean{PhyXMiniProblems.ProblemPhyXMini0281.lengthInCentimeters}
  \uses{def:physics:phyx-mini-0281:phyxminiproblems-problemphyxmini0281-lengthquantity, def:physics:phyx-mini-0281:phyxminiproblems-problemphyxmini0281-lengthreadout}
  Centimetre readout of a physical length
\end{definition}
\begin{proof}
  This is the declared model definition of length In Centimeters.
\end{proof}
\begin{definition}[mass In Kilograms]
  \label{def:physics:phyx-mini-0281:phyxminiproblems-problemphyxmini0281-massinkilograms}
  \lean{PhyXMiniProblems.ProblemPhyXMini0281.massInKilograms}
  \uses{def:physics:phyx-mini-0281:phyxminiproblems-problemphyxmini0281-massquantity}
  Kilogram readout of a physical mass
\end{definition}
\begin{proof}
  This is the declared model definition of mass In Kilograms.
\end{proof}
\begin{definition}[velocity In Meters Per Second]
  \label{def:physics:phyx-mini-0281:phyxminiproblems-problemphyxmini0281-velocityinmeterspersecond}
  \lean{PhyXMiniProblems.ProblemPhyXMini0281.velocityInMetersPerSecond}
  \uses{def:physics:phyx-mini-0281:phyxminiproblems-problemphyxmini0281-velocityquantity}
  Metre-per-second readout of a velocity magnitude
\end{definition}
\begin{proof}
  This is the declared model definition of velocity In Meters Per Second.
\end{proof}
\begin{definition}[acceleration In Meters Per Second Squared]
  \label{def:physics:phyx-mini-0281:phyxminiproblems-problemphyxmini0281-accelerationinmeterspersecondsquared}
  \lean{PhyXMiniProblems.ProblemPhyXMini0281.accelerationInMetersPerSecondSquared}
  \uses{def:physics:phyx-mini-0281:phyxminiproblems-problemphyxmini0281-accelerationquantity}
  Metre-per-second-squared readout of an acceleration magnitude
\end{definition}
\begin{proof}
  This is the declared model definition of acceleration In Meters Per Second Squared.
\end{proof}
\begin{definition}[spring Stiffness In Newtons Per Meter]
  \label{def:physics:phyx-mini-0281:phyxminiproblems-problemphyxmini0281-springstiffnessinnewtonspermeter}
  \lean{PhyXMiniProblems.ProblemPhyXMini0281.springStiffnessInNewtonsPerMeter}
  \uses{def:physics:phyx-mini-0281:phyxminiproblems-problemphyxmini0281-springstiffnessquantity}
  Newton-per-metre readout of a spring stiffness
\end{definition}
\begin{proof}
  This is the declared model definition of spring Stiffness In Newtons Per Meter.
\end{proof}
\begin{definition}[frequency In Hertz]
  \label{def:physics:phyx-mini-0281:phyxminiproblems-problemphyxmini0281-frequencyinhertz}
  \lean{PhyXMiniProblems.ProblemPhyXMini0281.frequencyInHertz}
  \uses{def:physics:phyx-mini-0281:phyxminiproblems-problemphyxmini0281-frequencyquantity}
  Hertz readout of a cyclic frequency
\end{definition}
\begin{proof}
  This is the declared model definition of frequency In Hertz.
\end{proof}
\begin{definition}[angular Frequency In Radians Per Second]
  \label{def:physics:phyx-mini-0281:phyxminiproblems-problemphyxmini0281-angularfrequencyinradianspersecond}
  \lean{PhyXMiniProblems.ProblemPhyXMini0281.angularFrequencyInRadiansPerSecond}
  \uses{def:physics:phyx-mini-0281:phyxminiproblems-problemphyxmini0281-angularfrequencyquantity}
  Radian-per-second readout of an angular frequency
\end{definition}
\begin{proof}
  This is the declared model definition of angular Frequency In Radians Per Second.
\end{proof}
\begin{definition}[Figure Panel]
  \label{def:physics:phyx-mini-0281:phyxminiproblems-problemphyxmini0281-figurepanel}
  \lean{PhyXMiniProblems.ProblemPhyXMini0281.FigurePanel}
  The three side-by-side configurations shown in the supplied image
\end{definition}
\begin{proof}
  This is the declared model definition of Figure Panel.
\end{proof}
\begin{definition}[Figure Label]
  \label{def:physics:phyx-mini-0281:phyxminiproblems-problemphyxmini0281-figurelabel}
  \lean{PhyXMiniProblems.ProblemPhyXMini0281.FigureLabel}
  The two mathematical labels printed in the image
\end{definition}
\begin{proof}
  This is the declared model definition of Figure Label.
\end{proof}
\begin{definition}[Figure Quantity Role]
  \label{def:physics:phyx-mini-0281:phyxminiproblems-problemphyxmini0281-figurequantityrole}
  \lean{PhyXMiniProblems.ProblemPhyXMini0281.FigureQuantityRole}
  The physical role of a printed figure label
\end{definition}
\begin{proof}
  This is the declared model definition of Figure Quantity Role.
\end{proof}
\begin{definition}[Spring Pendulum Figure]
  \label{def:physics:phyx-mini-0281:phyxminiproblems-problemphyxmini0281-springpendulumfigure}
  \lean{PhyXMiniProblems.ProblemPhyXMini0281.SpringPendulumFigure}
  \uses{def:physics:phyx-mini-0281:phyxminiproblems-problemphyxmini0281-figurepanel, def:physics:phyx-mini-0281:phyxminiproblems-problemphyxmini0281-figurelabel, def:physics:phyx-mini-0281:phyxminiproblems-problemphyxmini0281-figurequantityrole}
  Qualitative evidence visible directly in the supplied three-panel image
\end{definition}
\begin{proof}
  This is the declared model definition of Spring Pendulum Figure.
\end{proof}
\begin{definition}[Matches Primary Figure]
  \label{def:physics:phyx-mini-0281:phyxminiproblems-problemphyxmini0281-matchesprimaryfigure}
  \lean{PhyXMiniProblems.ProblemPhyXMini0281.MatchesPrimaryFigure}
  \uses{def:physics:phyx-mini-0281:phyxminiproblems-problemphyxmini0281-springpendulumfigure}
  Image readouts: the first two panels compare the same fixed spring before and after attaching the block, h marks their endpoint separation, and the third panel is a fixed-pivot string pendulum with a bob
\end{definition}
\begin{proof}
  This is the declared model definition of Matches Primary Figure.
\end{proof}
\begin{definition}[Motion Orientation]
  \label{def:physics:phyx-mini-0281:phyxminiproblems-problemphyxmini0281-motionorientation}
  \lean{PhyXMiniProblems.ProblemPhyXMini0281.MotionOrientation}
  Orientation of the block's oscillatory motion
\end{definition}
\begin{proof}
  This is the declared model definition of Motion Orientation.
\end{proof}
\begin{definition}[Oscillation Regime]
  \label{def:physics:phyx-mini-0281:phyxminiproblems-problemphyxmini0281-oscillationregime}
  \lean{PhyXMiniProblems.ProblemPhyXMini0281.OscillationRegime}
  Linearization regime used for either oscillator
\end{definition}
\begin{proof}
  This is the declared model definition of Oscillation Regime.
\end{proof}
\begin{definition}[Pendulum Model]
  \label{def:physics:phyx-mini-0281:phyxminiproblems-problemphyxmini0281-pendulummodel}
  \lean{PhyXMiniProblems.ProblemPhyXMini0281.PendulumModel}
  The ideal mechanical model requested for the pendulum
\end{definition}
\begin{proof}
  This is the declared model definition of Pendulum Model.
\end{proof}
\begin{definition}[Spring Pendulum Setup]
  \label{def:physics:phyx-mini-0281:phyxminiproblems-problemphyxmini0281-springpendulumsetup}
  \lean{PhyXMiniProblems.ProblemPhyXMini0281.SpringPendulumSetup}
  \uses{def:physics:phyx-mini-0281:phyxminiproblems-problemphyxmini0281-massquantity, def:physics:phyx-mini-0281:phyxminiproblems-problemphyxmini0281-lengthquantity, def:physics:phyx-mini-0281:phyxminiproblems-problemphyxmini0281-velocityquantity, def:physics:phyx-mini-0281:phyxminiproblems-problemphyxmini0281-accelerationquantity, def:physics:phyx-mini-0281:phyxminiproblems-problemphyxmini0281-springstiffnessquantity, def:physics:phyx-mini-0281:phyxminiproblems-problemphyxmini0281-frequencyquantity, def:physics:phyx-mini-0281:phyxminiproblems-problemphyxmini0281-angularfrequencyquantity, def:physics:phyx-mini-0281:phyxminiproblems-problemphyxmini0281-springpendulumfigure, def:physics:phyx-mini-0281:phyxminiproblems-problemphyxmini0281-motionorientation, def:physics:phyx-mini-0281:phyxminiproblems-problemphyxmini0281-oscillationregime, def:physics:phyx-mini-0281:phyxminiproblems-problemphyxmini0281-pendulummodel}
  All independent physical quantities used by the spring and pendulum models. The pendulum length and both frequencies are unconstrained outputs here. In particular, neither the extension h nor its 2 cm readout is assigned to the pendulum length by this structure
\end{definition}
\begin{proof}
  This is the declared model definition of Spring Pendulum Setup.
\end{proof}
\begin{definition}[Matches Problem Description]
  \label{def:physics:phyx-mini-0281:phyxminiproblems-problemphyxmini0281-matchesproblemdescription}
  \lean{PhyXMiniProblems.ProblemPhyXMini0281.MatchesProblemDescription}
  \uses{def:physics:phyx-mini-0281:phyxminiproblems-problemphyxmini0281-lengthinmeters, def:physics:phyx-mini-0281:phyxminiproblems-problemphyxmini0281-lengthincentimeters, def:physics:phyx-mini-0281:phyxminiproblems-problemphyxmini0281-velocityinmeterspersecond, def:physics:phyx-mini-0281:phyxminiproblems-problemphyxmini0281-springpendulumsetup}
  Numerical and qualitative data stated in the prose. “A short distance” is recorded as a positive pull in the small-vertical-oscillation regime; release means zero initial velocity. The standard small-angle regime is made explicit for the simple-pendulum formula. No requested pendulum length occurs here
\end{definition}
\begin{proof}
  This is the declared model definition of Matches Problem Description.
\end{proof}
\begin{definition}[Has Physical Parameters]
  \label{def:physics:phyx-mini-0281:phyxminiproblems-problemphyxmini0281-hasphysicalparameters}
  \lean{PhyXMiniProblems.ProblemPhyXMini0281.HasPhysicalParameters}
  \uses{def:physics:phyx-mini-0281:phyxminiproblems-problemphyxmini0281-lengthinmeters, def:physics:phyx-mini-0281:phyxminiproblems-problemphyxmini0281-massinkilograms, def:physics:phyx-mini-0281:phyxminiproblems-problemphyxmini0281-accelerationinmeterspersecondsquared, def:physics:phyx-mini-0281:phyxminiproblems-problemphyxmini0281-springstiffnessinnewtonspermeter, def:physics:phyx-mini-0281:phyxminiproblems-problemphyxmini0281-frequencyinhertz, def:physics:phyx-mini-0281:phyxminiproblems-problemphyxmini0281-angularfrequencyinradianspersecond, def:physics:phyx-mini-0281:phyxminiproblems-problemphyxmini0281-springpendulumsetup}
  Positivity and nondegeneracy conditions for the physical branch
\end{definition}
\begin{proof}
  This is the declared model definition of Has Physical Parameters.
\end{proof}
\begin{definition}[spring Oscillator SI]
  \label{def:physics:phyx-mini-0281:phyxminiproblems-problemphyxmini0281-springoscillatorsi}
  \lean{PhyXMiniProblems.ProblemPhyXMini0281.springOscillatorSI}
  \uses{def:physics:phyx-mini-0281:phyxminiproblems-problemphyxmini0281-massinkilograms, def:physics:phyx-mini-0281:phyxminiproblems-problemphyxmini0281-springstiffnessinnewtonspermeter, def:physics:phyx-mini-0281:phyxminiproblems-problemphyxmini0281-springpendulumsetup, def:physics:phyx-mini-0281:phyxminiproblems-problemphyxmini0281-hasphysicalparameters}
  The loaded spring viewed in coherent SI readouts as Physlib's classical harmonic oscillator, whose angular frequency is sqrt (k / m)
\end{definition}
\begin{proof}
  This is the declared model definition of spring Oscillator SI.
\end{proof}
\begin{definition}[pendulum Oscillator SI]
  \label{def:physics:phyx-mini-0281:phyxminiproblems-problemphyxmini0281-pendulumoscillatorsi}
  \lean{PhyXMiniProblems.ProblemPhyXMini0281.pendulumOscillatorSI}
  \uses{def:physics:phyx-mini-0281:phyxminiproblems-problemphyxmini0281-lengthinmeters, def:physics:phyx-mini-0281:phyxminiproblems-problemphyxmini0281-massinkilograms, def:physics:phyx-mini-0281:phyxminiproblems-problemphyxmini0281-accelerationinmeterspersecondsquared, def:physics:phyx-mini-0281:phyxminiproblems-problemphyxmini0281-springpendulumsetup, def:physics:phyx-mini-0281:phyxminiproblems-problemphyxmini0281-hasphysicalparameters}
  The small-angle simple pendulum viewed as a generalized harmonic oscillator. For angular displacement, generalized inertia is m L² and the gravitational restoring coefficient is m g L, so their ratio is g / L
\end{definition}
\begin{proof}
  This is the declared model definition of pendulum Oscillator SI.
\end{proof}
\begin{definition}[Satisfies Spring And Pendulum Laws]
  \label{def:physics:phyx-mini-0281:phyxminiproblems-problemphyxmini0281-satisfiesspringandpendulumlaws}
  \lean{PhyXMiniProblems.ProblemPhyXMini0281.SatisfiesSpringAndPendulumLaws}
  \uses{def:physics:phyx-mini-0281:phyxminiproblems-problemphyxmini0281-lengthinmeters, def:physics:phyx-mini-0281:phyxminiproblems-problemphyxmini0281-massinkilograms, def:physics:phyx-mini-0281:phyxminiproblems-problemphyxmini0281-accelerationinmeterspersecondsquared, def:physics:phyx-mini-0281:phyxminiproblems-problemphyxmini0281-springstiffnessinnewtonspermeter, def:physics:phyx-mini-0281:phyxminiproblems-problemphyxmini0281-frequencyinhertz, def:physics:phyx-mini-0281:phyxminiproblems-problemphyxmini0281-angularfrequencyinradianspersecond, def:physics:phyx-mini-0281:phyxminiproblems-problemphyxmini0281-springpendulumsetup, def:physics:phyx-mini-0281:phyxminiproblems-problemphyxmini0281-hasphysicalparameters, def:physics:phyx-mini-0281:phyxminiproblems-problemphyxmini0281-springoscillatorsi, def:physics:phyx-mini-0281:phyxminiproblems-problemphyxmini0281-pendulumoscillatorsi}
  Static balance and the two general small-oscillation laws. The static law is k h = m g. Physlib supplies ω = sqrt (k / m) for the spring and for the generalized pendulum oscillator, and the last two fields state ω = 2πf. None of these fields gives a numerical pendulum length or asserts that the pendulum length equals h
\end{definition}
\begin{proof}
  This is the declared model definition of Satisfies Spring And Pendulum Laws.
\end{proof}
\begin{definition}[Pendulum Matches Spring Frequency]
  \label{def:physics:phyx-mini-0281:phyxminiproblems-problemphyxmini0281-pendulummatchesspringfrequency}
  \lean{PhyXMiniProblems.ProblemPhyXMini0281.PendulumMatchesSpringFrequency}
  \uses{def:physics:phyx-mini-0281:phyxminiproblems-problemphyxmini0281-frequencyinhertz, def:physics:phyx-mini-0281:phyxminiproblems-problemphyxmini0281-springpendulumsetup}
  The design condition posed by the question: the simple pendulum swings with the same cyclic frequency as the spring--block oscillator. This constrains a frequency, not the requested length
\end{definition}
\begin{proof}
  This is the declared model definition of Pendulum Matches Spring Frequency.
\end{proof}
\begin{definition}[Answer Choice]
  \label{def:physics:phyx-mini-0281:phyxminiproblems-problemphyxmini0281-answerchoice}
  \lean{PhyXMiniProblems.ProblemPhyXMini0281.AnswerChoice}
  Labels of the four answer choices printed in the source
\end{definition}
\begin{proof}
  This is the declared model definition of Answer Choice.
\end{proof}
\begin{definition}[length In Centimeters]
  \label{def:physics:phyx-mini-0281:phyxminiproblems-problemphyxmini0281-answerchoice-lengthincentimeters}
  \lean{PhyXMiniProblems.ProblemPhyXMini0281.AnswerChoice.lengthInCentimeters}
  \uses{def:physics:phyx-mini-0281:phyxminiproblems-problemphyxmini0281-lengthincentimeters, def:physics:phyx-mini-0281:phyxminiproblems-problemphyxmini0281-answerchoice}
  Pendulum length in centimetres printed beside an answer label
\end{definition}
\begin{proof}
  This is the declared model definition of length In Centimeters.
\end{proof}
\begin{definition}[recorded Answer Choice]
  \label{def:physics:phyx-mini-0281:phyxminiproblems-problemphyxmini0281-recordedanswerchoice}
  \lean{PhyXMiniProblems.ProblemPhyXMini0281.recordedAnswerChoice}
  \uses{def:physics:phyx-mini-0281:phyxminiproblems-problemphyxmini0281-answerchoice}
  The answer label recorded in the dataset metadata
\end{definition}
\begin{proof}
  This is the declared model definition of recorded Answer Choice.
\end{proof}
\begin{definition}[Matches Answer Choice]
  \label{def:physics:phyx-mini-0281:phyxminiproblems-problemphyxmini0281-matchesanswerchoice}
  \lean{PhyXMiniProblems.ProblemPhyXMini0281.MatchesAnswerChoice}
  \uses{def:physics:phyx-mini-0281:phyxminiproblems-problemphyxmini0281-lengthincentimeters, def:physics:phyx-mini-0281:phyxminiproblems-problemphyxmini0281-springpendulumsetup, def:physics:phyx-mini-0281:phyxminiproblems-problemphyxmini0281-answerchoice}
  A displayed choice reports the physical pendulum length exactly
\end{definition}
\begin{proof}
  This is the declared model definition of Matches Answer Choice.
\end{proof}
\begin{lemma}[matching Frequency pendulum Length eq equilibrium Extension]
  \label{lem:physics:phyx-mini-0281:phyxminiproblems-problemphyxmini0281-matchingfrequency-pendulumlength-eq-equilibriumextension}
  \lean{PhyXMiniProblems.ProblemPhyXMini0281.matchingFrequency_pendulumLength_eq_equilibriumExtension}
  \uses{def:physics:phyx-mini-0281:phyxminiproblems-problemphyxmini0281-springpendulumsetup, def:physics:phyx-mini-0281:phyxminiproblems-problemphyxmini0281-hasphysicalparameters, def:physics:phyx-mini-0281:phyxminiproblems-problemphyxmini0281-satisfiesspringandpendulumlaws, def:physics:phyx-mini-0281:phyxminiproblems-problemphyxmini0281-pendulummatchesspringfrequency}
  Any positive simple pendulum satisfying the stated laws and frequency-match condition has physical length equal to the spring's equilibrium extension. This is an intermediate conclusion, not an assumption
\end{lemma}
\begin{proof}
  The conclusion follows from the governing relations and figure readouts named in the statement of matching Frequency pendulum Length eq equilibrium Extension.
\end{proof}
% --- Archon declaration topology end ---
% --- Archon physics formalization source end ---
